\section{Aerodynamic Drag (CdA) Fitting}
\label{sec:cda-fitting}

\texttt{CdAFitter} inverts the equation of motion of
\cref{sec:simulation-model,sec:cornering-model}: instead of integrating forward
from a known \CdA{} to predict $v(t)$, it takes a measured $\bigl(t, v, P\bigr)$
trace from a rider's \texttt{.fit} file and solves for the one remaining
unknown, \CdA{}.

\subsection{Problem statement}
\label{sec:cda-problem}

Given a selected window of recorded time, speed, and power samples, the rider's
acceleration can be computed directly by differentiating the recorded speed
trace, leaving \cref{eq:newton-second-law} (extended with the cornering forces
of \cref{sec:cornering-model}) as a single linear equation in the one unknown
$\CdA$, since every other quantity in the force balance -- mass, rolling
resistance, propulsive force, and the cornering corrections -- can be computed
directly from the recorded data and the rider's/environment's known parameters.

\subsection{Acceleration from recorded data}
\label{sec:cda-dvdt}

Recorded samples are filtered to those with $v > V_{\mathrm{eps}}$, and the
acceleration $dv/dt$ is estimated by a central finite difference
(\texttt{numpy.gradient}) on the raw, non-uniformly spaced recording timestamps,
with no smoothing applied to the speed trace beforehand. Distance traveled
within the selected window is obtained by cumulative trapezoidal integration of
the recorded speed, rather than from a separate odometer or GPS-distance field.
Both choices are revisited as limitations in \cref{sec:assumptions}.

\subsection{Rearranged force balance}
\label{sec:cda-rearranged}

Rearranging \cref{eq:newton-second-law} (with $F_{ad}$ from
\cref{eq:aero-drag} and the cornering-corrected rolling resistance and energy
terms of \cref{sec:cornering-model}) to isolate the drag force gives
\begin{equation}
    F_{ad} = F_p(1-\eta) + F_{rr}(\theta) + m\,a_{\mathrm{energy}} - m\,\frac{dv}{dt},
    \label{eq:cda-force-balance}
\end{equation}
where the propulsive force here uses the recorded power directly,
$F_p = P(t)/v$, with no traction cap applied (unlike the forward model's
\cref{eq:propulsive-force}), since $P(t)$ is measured rather than a target to be
achieved. Since $F_{ad} = \CdA \cdot \bigl(\tfrac{1}{2}\rho v^2\bigr)$ by
\cref{eq:aero-drag}, writing
\begin{equation}
    x \equiv \tfrac{1}{2}\rho v^2, \qquad
    y \equiv F_p(1-\eta) + F_{rr}(\theta) + m\,a_{\mathrm{energy}} - m\,\frac{dv}{dt},
    \label{eq:cda-xy}
\end{equation}
reduces the fit to a single linear regression, $y = \CdA \cdot x$, evaluated
across every sample in the selected window.

\subsection{Closed-form least-squares fit}
\label{sec:cda-ols}

Because \cref{eq:aero-drag} requires $F_{ad} = 0$ at $v = 0$, the correct
regression has no intercept term: the fit is constrained to pass through the
origin, not merely close to it. The ordinary-least-squares solution to
$y = \CdA \cdot x$ with a fixed zero intercept has the closed form
\begin{equation}
    \CdA = \frac{\sum_i x_i y_i}{\sum_i x_i^2},
    \label{eq:cda-closed-form}
\end{equation}
computed directly with no iterative solver:
\begin{lstlisting}
cda = np.sum(x * y) / np.sum(x ** 2)
residual = np.sum((y - x * cda) ** 2)
\end{lstlisting}
The fit is rejected if $\CdA \leq 0$ (a non-physical result) or if the selected
window has too little speed variation for $\sum_i x_i^2$ to meaningfully
constrain the regression.

\subsection{Lap-phase search}
\label{sec:cda-s0-search}

The track's curvature $\kappa(s)$ is periodic in position $s$
(\cref{sec:track-geometry}), but the phase offset $s_0$ at which the recorded
window began -- where on the track, modulo one straight-and-corner period, the
rider was when recording started -- is not directly known from the \texttt{.fit}
data. \texttt{CdAFitter} searches for $s_0$ in two stages:

\begin{enumerate}
    \item \textbf{Coarse grid search}: the fit of
    \cref{eq:cda-closed-form} is evaluated at $150$ candidate offsets spaced
    evenly over \emph{half} the track's straight-and-corner period. Only half
    the period needs to be searched because the track's two straights and two
    corners are geometrically symmetric within a full lap.

    \item \textbf{Refinement}: the best coarse candidate is refined with
    \url{scipy.optimize.minimize_scalar} using the bounded Brent method,
    searching within one grid spacing on either side of the coarse minimum.
\end{enumerate}
At each candidate $s_0$, the scoring function is the sum of squared residuals of
the fit in \cref{eq:cda-closed-form},
$\sum_i \bigl(y_i - x_i \CdA(s_0)\bigr)^2$, so $s_0$ and $\CdA$ are fit jointly:
$\CdA$ is solved in closed form for each candidate $s_0$, and $s_0$ itself is
chosen to minimize the resulting residual. If no track geometry is configured,
$s_0$ is not searched and the fit is evaluated directly with $\kappa \equiv 0$
throughout, in which case the cornering terms of
\cref{eq:cda-force-balance} vanish identically.

\subsection{Sample requirements}
\label{sec:cda-requirements}

The fit requires at least two samples with $v > V_{\mathrm{eps}}$ after
filtering; a selection window with fewer moving samples raises an error rather
than returning a degenerate fit.
